\raggedright

\section{Overview of the Project}

This project presents the design and implementation of TAC, an identity-centric case management system developed as a secure infrastructure-backed platform for handling investigative case and subject records. The system was designed to demonstrate how a modern case management environment can be protected through a combination of Zero Trust principles, role-based access control, federated authentication, network segmentation, audit logging, and monitored infrastructure.

The project is not limited to the development of a web application. Instead, it combines both software engineering and infrastructure engineering to create a realistic enterprise-style environment. The final system includes a Next.js web application, a Microsoft SQL Server database, Active Directory Domain Services, Active Directory Federation Services, pfSense firewalling, VLAN-based segmentation, SSL/TLS certificates, Wazuh monitoring, Kasm virtual access, and MinIO-style audit storage concepts. These components work together to support the wider project goal of demonstrating how sensitive case information can be protected using strong identity and infrastructure controls.

TAC was developed as a case and subject registry platform. The application allows authorised users to create, view, update, delete, import, and analyse case-related records. The system also supports audit visibility, allowing user activity to be recorded and reviewed. From a user perspective, the platform appears as a secure internal web application. From a technical perspective, it is supported by multiple infrastructure layers that control who can access the system, from where access is allowed, how data is stored, and how activity is monitored.

The project was developed over two semesters. Semester One focused primarily on research, design, planning, and architectural validation. Semester Two focused on implementation, integration, testing, refinement, and final validation. This report documents the completed Semester Two implementation and explains the project from initial concept through to final system design and testing.

% Suggested figure:
% Figure 1.1: Screenshot of the TAC login page or main dashboard.
% Caption: TAC secure case management dashboard.

\section{Project Context}

Investigative and intelligence led environments routinely process information that may be sensitive, private, operationally significant, or legally relevant. This information can include case details, subject records, internal notes, audit trails, operational classifications, and links between people, cases, and events. In such environments, the security of the system cannot depend only on a traditional network boundary. A user being inside a network does not automatically mean that the user should have access to all records.

Many older systems are based on the assumption that internal networks are trusted and external networks are untrusted. This model is increasingly unsuitable for modern systems because users may work remotely, access systems through VPNs, use cloud services, or authenticate across multiple domains and applications. In addition, threats may come from compromised accounts, misconfigured services, excessive administrator privileges, weak monitoring, or poor separation between user, server, and management networks.

TAC was developed in response to this problem. The project demonstrates how a secure case management platform can be designed around identity rather than location. In this model, access is controlled through verified user identity, group membership, assigned role, network path, and system policy. The platform is intended to show how case data can be protected using multiple layers of security rather than relying on a single control.

The wider project context also includes digital identity, data protection, and accountability. A secure case management system must be able to answer important questions such as who accessed a record, what action was performed, when the action occurred, and whether the user was authorised to perform it. These questions are especially important where personal or sensitive information is involved. Therefore, the system was designed with authentication, authorisation, auditability, and monitoring as core requirements rather than optional additions.

% Suggested figure:
% Figure 1.2: High-level context diagram showing users, identity provider, web app, SQL database, and monitoring.
% Caption: High-level context of the TAC platform.

\section{Motivation and Problem Statement}

The motivation for this project came from the need to design a realistic secure information system that brings together infrastructure security and application security. A case management system is a suitable focus because it naturally involves sensitive data, different user roles, audit requirements, and strict access control. In a real organisation, not every user should have the same level of access. For example, an investigator may require access to create and update case records, a supervisor may require wider visibility and approval rights, while an intelligence reader may only require read only access to selected information.

The central problem addressed by this project is how to design and implement a secure case management platform where access is based on identity, role, and least privilege rather than simple network presence. A secondary problem is how to support this application with realistic infrastructure services, including Active Directory, federation, SQL storage, firewall controls, monitoring, and controlled remote access.

A traditional web application by itself would not fully address these requirements. If the application stored data locally, had no role based access model, lacked audit logging, or was hosted on a flat network, it would not reflect the security needs of a real investigative environment. For this reason, TAC was developed as a full stack and infrastructure backed project. The application layer and infrastructure layer were designed to support each other.

The problem can therefore be summarised as follows: sensitive case and subject information requires a secure platform that can control access, separate infrastructure zones, authenticate users through a central identity provider, store data reliably, monitor system activity, and provide evidence that the implemented controls function as intended. This project addresses that problem by implementing TAC as a proof of concept secure case management platform.

% Suggested figure:
% Figure 1.3: Problem diagram comparing a flat/traditional access model against an identity-centric Zero Trust model.
% Caption: Traditional perimeter trust compared with identity-centric access.

\section{Aim of the Project}

The aim of this project is to design, implement, and validate a secure identity-centric case management system supported by enterprise style infrastructure, which mirrors the confidentiality steps taken  to reflect law enforcement and Private Investigation systems. The project aims to demonstrate how Zero Trust principles can be applied in practice to a case management environment by combining identity management, role based access control, network segmentation, secure database storage, audit logging, and infrastructure monitoring.

The system was developed to show that secure application design should not be separated from infrastructure design. The web application provides the user facing case management functionality, while the underlying infrastructure provides the controls that support authentication, authorization, segmentation, monitoring, and secure communication. By combining these areas, the project demonstrates a more complete security model than a standalone prototype application.

A further aim is to provide a realistic final year project implementation that reflects the type of systems used in enterprise and security focused environments. The project therefore includes both practical implementation and technical validation. The final system is evaluated not only by whether the application functions correctly, but also by whether the supporting infrastructure behaves as intended.

% Suggested figure:
% Figure 1.4: System aim diagram showing application security and infrastructure security joining together.
% Caption: TAC project aim combining application and infrastructure security.

\section{Project Objectives}

The project objectives were defined to cover both the software system and the supporting infrastructure. The first objective was to design and build a secure web based case management application that allows authorised users to manage case and subject records through a clear and usable interface. This included functionality for creating, viewing, editing, deleting, importing, analysing records and most importantly view the audit trails and logs of all activities that happen inside the case management platform.

The second objective was to replace early prototype storage with a persistent SQL-backed data model. Microsoft SQL Server was selected to provide structured relational storage for case, subject, and audit related data. This allowed the application to move beyond local mock data and operate as a more realistic enterprise system.

The third objective was to implement identity-based access using Active Directory and Active Directory Federation Services. This objective focused on using centralised identity, across two different domain controllers which helps split domain user responsibility based on roles, security groups, claims, and role mapping to control access to the application.

The fourth objective was to design a segmented infrastructure using pfSense and VLANs. This was required to separate users, servers, management systems, VPN access, and security monitoring. The purpose of this objective was to reduce unnecessary network exposure and ensure that systems such as SQL Server were not directly accessible to ordinary users.

The fifth objective was to implement monitoring and auditability. This included application-level audit logging and infrastructure level monitoring through Wazuh. The purpose was to ensure that both user activity and system events could be observed and reviewed such as users logging in and out of management and security servers.

The sixth objective was to support controlled access through technologies such as a VPN. This reflected the requirement for secure remote or field based access without exposing sensitive systems directly.

The final objective was to validate the completed system through testing. Each major feature and infrastructure control was tested to confirm that the system behaved as expected. Validation included functional testing, role-based access testing, SQL connectivity testing, firewall and VLAN testing, SSL testing, monitoring validation, and synthetic data testing.

% Suggested figure:
% Figure 1.5: Objectives table or visual checklist showing completed project objectives.
% Caption: Summary of TAC project objectives.

\section{Scope of the Final Implementation}

The final implementation includes both the TAC web application and the infrastructure required to support it. The application scope includes the case registry, subject registry, case and subject detail views, create and edit forms, delete functionality, bulk import functionality, analytics, audit log visibility, and a structured dashboard interface. The system was designed to be usable by different roles within an investigative-style organisation.

The infrastructure scope includes a segmented virtual lab environment using pfSense, VLANs, Proxmox-hosted systems, Windows Server infrastructure, Active Directory, ADFS, SQL Server, Linux-based web hosting, Wazuh monitoring. SSL/TLS certificates were also configured to support encrypted access to internal services that were self signed as its managed as a private intranet.

The scope also includes the use of synthetic data for testing and demonstration. No real case data, personal investigation data, suspect data, victim data, or operational law enforcement data was used. Synthetic records were generated to test the bulk import process, analytics, search, audit logging, and general application workflows. This approach allowed the system to be demonstrated realistically while avoiding ethical and data protection risks.

Some elements were implemented as proof of concept or design-supported components rather than full production systems. For example, disaster recovery planning, MinIO-style append-only audit storage, and business continuity planning were considered as part of the wider architecture but are not presented as fully production hardened enterprise deployments. They are included to show how the platform could be extended into a more mature operational environment.

This scope ensures that the project remains achievable within the final year project timeframe while still demonstrating a broad range of technical skills across software development, systems administration, identity management, network security, and infrastructure validation.

% Suggested figure:
% Figure 1.6: Scope boundary diagram showing what is inside and outside the final implementation.
% Caption: Scope of the TAC final implementation.

\section{Structure of the Report}

This report is structured to explain the project from initial concept through to completed implementation and validation. The Introduction chapter provides an overview of the project, its context, motivation, aims, objectives, and final scope.

The Requirements Analysis chapter describes the business, functional, non-functional, security, identity, infrastructure, data management, audit, and data protection requirements that shaped the system. This chapter explains what the platform needed to achieve before discussing the technical design.

The High Level Design chapter presents the first-cut design decisions and explains the technology choices made for the project. It discusses why specific technologies such as Active Directory, ADFS, SQL Server, Next.js, pfSense, Wazuh and MinIO were selected.

The Systems Architecture chapter provides the detailed design of the completed system. It covers the infrastructure architecture, network design, VLAN layout, firewall controls, identity architecture, federation design, application architecture, database architecture, monitoring design, and disaster recovery considerations.

The Prototyping and Wireframes chapter explains how the system evolved through mock data, interface prototyping, infrastructure prototyping, and iterative refinement. It also discusses the role of wireframes and supporting AI tools during early design exploration for the UI interface and bulk mock data generated.

The Implementation chapter documents how the major components were built and configured. This includes the network, servers, Active Directory, ADFS, SQL Server, SSL certificates, Next.js application, bulk import, analytics, audit logging, RBAC, Wazuh, and MinIO-related components.

The Validation and Testing chapter demonstrates that the implemented features and infrastructure controls were tested. This chapter includes functional testing, security testing, role testing, SQL connectivity testing, firewall testing, VLAN testing, SSL testing, monitoring validation, Selenium testing, and synthetic data validation.

The Ethical, Legal and Data Protection Considerations chapter explains how the project avoided the use of real personal data, how synthetic data was used, and how GDPR principles such as minimisation, access control, accountability, and auditability influenced the design.

The Project Management and Version Control chapter outlines the development methodology, GitHub usage, branching strategy, commit history, task planning, and change management. The report then concludes with reflection, future work, and a final evaluation of the completed project.

% Suggested figure:
% Figure 1.7: Report structure diagram showing the chapters as a flow from requirements to validation.
% Caption: Structure of the final report.
